% Generated by docs/generate_user_guide.py. Do not edit directly.
% Source: utilities/descriptions.yaml: analysis_parameters

\begin{longtable}{@{}>{\raggedright\arraybackslash}p{0.22\textwidth}>{\raggedright\arraybackslash}p{0.22\textwidth}>{\raggedright\arraybackslash}p{0.48\textwidth}@{}}
\toprule
\textbf{Parameter} & \textbf{Internal key} & \textbf{Description} \\
\midrule
\endfirsthead
\toprule
\textbf{Parameter} & \textbf{Internal key} & \textbf{Description} \\
\midrule
\endhead
Cross-validation scheme & cv\_scheme & Selects the splitting method used to estimate held-out performance for fitted parametric models. Leave-one-timepoint-out is the default. \\
CV splits/folds/block size & cv\_n\_splits & For Shuffle split, the number of repeated random train/test splits. For K-fold, the number of folds. For Leave-contiguous-block-out, the number of adjacent time points held out per split. The setup page does not enforce a maximum block size because the uploaded data are not read until Analyze is selected; larger blocks may leave too few observations for stable fitting. This setting is not used for Leave-one-timepoint-out or None. \\
Shuffle split test size & cv\_test\_size & The fraction of observations held out for testing in each Shuffle split repeat. Smaller values leave more observations for fitting; larger values make the holdout test more demanding but leave fewer observations for training. \\
CV random seed & cv\_random\_state & Seed used to make randomized cross-validation splits reproducible. This applies to Shuffle split. \\
\bottomrule
\end{longtable}
